\section{Canonical quasicoherent affine-cover effacement}

\begin{lemma}[Finite products of quasicoherent modules]
\label{lem:scheme-quasicoherent-finite-product}
\lean{EnvelopeProperty.SchemeQuasicoherentAffineEffacement.isQuasicoherent_product}
\leanok
Let $X$ be a scheme and $(M_i)_{i\in I}$ a finite family of actual
quasicoherent $\mathcal O_X$-modules, with $I$ in the same universe as
$X$. The categorical product $\prod_i M_i$ is quasicoherent.
\end{lemma}
\begin{proof}
Restrict the actual product to each affine open of $X$. Native open
restriction preserves limits, hence identifies this restriction with
the product of the restrictions of the original $M_i$. Quasicoherent
modules on an affine scheme are closed under finite products. Transport
this property along that comparison, then use the affine-open cover
criterion for quasicoherence. No affineness hypothesis on $X$ is used.
\end{proof}

\begin{definition}[The original cover augmentation and its cokernel]
\label{def:scheme-canonical-affine-effacement}
\lean{EnvelopeProperty.SchemeQuasicoherentAffineEffacement.E,EnvelopeProperty.SchemeQuasicoherentAffineEffacement.augmentation,EnvelopeProperty.SchemeQuasicoherentAffineEffacement.quotient,EnvelopeProperty.SchemeQuasicoherentAffineEffacement.π,EnvelopeProperty.SchemeQuasicoherentAffineEffacement.sequence}
\leanok
For an original family $U:I\to X.\mathrm{Opens}$, put $j_i:U_i\hookrightarrow X$
and define
\[
 E_U(M)=\prod_{i\in I}(j_i)_*j_i^*M,\qquad
 a_M:M\longrightarrow E_U(M),\qquad Q_U(M)=\operatorname{coker}(a_M).
\]
Here $a_M$ is the native map \texttt{M.toCoverSheaf U}: its $i$th
component is the original restriction--pushforward adjunction unit.
Write $\pi_M:E_U(M)\to Q_U(M)$ for the actual cokernel projection.
The sequence is the actual short complex $(a_M,\pi_M)$.
These definitions require neither quasicoherence nor geometric hypotheses.
\end{definition}

\begin{theorem}[A canonical short exact sequence inside quasicoherent modules]
\label{thm:scheme-affine-effacement-short-exact}
\lean{EnvelopeProperty.SchemeQuasicoherentAffineEffacement.augmentation_mono,EnvelopeProperty.SchemeQuasicoherentAffineEffacement.E_isQuasicoherent,EnvelopeProperty.SchemeQuasicoherentAffineEffacement.quotient_isQuasicoherent,EnvelopeProperty.SchemeQuasicoherentAffineEffacement.sequence_shortExact,EnvelopeProperty.SchemeQuasicoherentAffineEffacement.sheafSequence_shortExact}
\leanok
\uses{def:scheme-canonical-affine-effacement,lem:scheme-quasicoherent-finite-product}
Suppose $X$ is separated, $I$ finite, the original $U_i$ are affine
and cover $X$, and the original $M$ is quasicoherent. Then
\[
 0\longrightarrow M\xrightarrow{a_M}E_U(M)
 \xrightarrow{\pi_M}Q_U(M)\longrightarrow0
\]
is short exact, both $E_U(M)$ and $Q_U(M)$ are quasicoherent, and
the same actual sequence of underlying additive sheaves is short exact.
\end{theorem}
\begin{proof}
The native cover-augmentation theorem makes this exact $a_M$ monic.
Since $X$ is separated, each affine-open inclusion $j_i$ is affine:
its composite with $X\to\operatorname{Spec}\mathbb Z$ is affine,
and affine morphisms descend through separated postcomposition.
Restriction preserves quasicoherence and affine pushforward preserves
quasicoherence. The finite-product lemma therefore applies to $E_U(M)$.
The native cokernel closure theorem gives quasicoherence of $Q_U(M)$.
The cokernel complex of a monomorphism is short exact. Finally the
actual module-to-additive-sheaf functor is exact, so its image of this
same short complex is short exact.
\end{proof}

\begin{theorem}[Actual positive cohomology of the effacing module vanishes]
\label{thm:scheme-affine-effacement-acyclic}
\lean{EnvelopeProperty.SchemeQuasicoherentAffineEffacement.E_subsingleton_H}
\leanok
\uses{thm:scheme-affine-effacement-short-exact,thm:affine-pushforward-cohomology-vanishing}
Under the preceding original hypotheses, for every $q\ge0$,
the actual small sheaf-cohomology group $H^{q+1}(X,E_U(M))$ is zero.
Neither a cohomology-vanishing certificate nor a larger-universe
cohomology comparison is an input.
\end{theorem}
\begin{proof}
The restriction $j_i^*M$ is quasicoherent on the original affine $U_i$,
so all its positive cohomology vanishes by affine acyclicity.
The genuine affine-pushforward comparison transfers this vanishing to
$(j_i)_*j_i^*M$ on $X$. An order on the finite original index set is
chosen internally solely to apply that comparison; no order is part
of the statement. The additive-sheaf functor preserves the actual
finite product defining $E_U(M)$. Actual small sheaf cohomology also
preserves finite products. Thus the cohomology in question is
isomorphic to a finite product of zero groups, and is zero.
The empty finite index set is included whenever it covers $X$.
\end{proof}

\begin{theorem}[Whole functoriality for the original augmentation and quotient]
\label{thm:scheme-affine-effacement-naturality}
\lean{EnvelopeProperty.SchemeQuasicoherentAffineEffacement.EMap,EnvelopeProperty.SchemeQuasicoherentAffineEffacement.EMap_π,EnvelopeProperty.SchemeQuasicoherentAffineEffacement.EMap_id,EnvelopeProperty.SchemeQuasicoherentAffineEffacement.EMap_comp,EnvelopeProperty.SchemeQuasicoherentAffineEffacement.EFunctor,EnvelopeProperty.SchemeQuasicoherentAffineEffacement.augmentation_naturality,EnvelopeProperty.SchemeQuasicoherentAffineEffacement.quotientMap,EnvelopeProperty.SchemeQuasicoherentAffineEffacement.π_naturality,EnvelopeProperty.SchemeQuasicoherentAffineEffacement.quotientMap_id,EnvelopeProperty.SchemeQuasicoherentAffineEffacement.quotientMap_comp,EnvelopeProperty.SchemeQuasicoherentAffineEffacement.quotientFunctor}
\leanok
\uses{def:scheme-canonical-affine-effacement}
For every actual original morphism $f:M\to N$, the product of the
restriction--pushforward images of $f$ defines $E_U(f)$ and satisfies
\[
 f\,a_N=a_M E_U(f).
\]
The induced actual cokernel morphism $Q_U(f)$ satisfies
\[
 \pi_M Q_U(f)=E_U(f)\pi_N.
\]
Both constructions preserve identities and composition, so they give
endofunctors on the original category $X.\mathrm{Modules}$. All these
are whole sheaf morphism identities; no matching hypothesis is supplied.
\end{theorem}
\begin{proof}
Define $E_U(f)$ by the actual product universal property. Its composite
with each product projection is the corresponding restricted and
pushforward map. Naturality of each native adjunction unit gives the
first square after these projections, hence as a whole identity.
The actual cokernel universal property then defines $Q_U(f)$ and gives
the second square. Product projections detect the identity and
composition laws for $E_U$; epimorphic cokernel projections detect
the same laws for $Q_U$. No quasicoherence or separation is needed
for this functoriality.
\end{proof}
